% The blueprint content, shared by the web and print versions.
%
% Node conventions (leanblueprint):
%   \lean{Name}   ties the node to a Lean declaration; `leanblueprint checkdecls`
%                 verifies the name exists in the built QuantumBook library.
%   \leanok       marks the statement/proof as fully formalized (green in graph).
%   \uses{labels} records dependencies (edges in the graph).
%   A node with neither \lean nor \leanok is a FRONTIER node (prompt.md invariant
%   2): a result not yet available in Mathlib, tracked but not proven.

\chapter{Part 0 -- Substrate}

This chapter fixes the finite-dimensional linear algebra on which quantum
mechanics is built, and its organizing result: the spectral theorem for
self-adjoint operators. Everything here is fully verified against Mathlib.

\section{The finite-dimensional spectral theorem}

\begin{theorem}[Observables have real spectra]
  \label{thm:selfadjoint-real}
  \lean{QuantumBook.Foundations.selfAdjoint_eigenvalue_isReal}
  \leanok
  Let $T$ be a self-adjoint (symmetric) operator on a complex inner product
  space $E$. Every eigenvalue $\mu$ of $T$ is real: $\operatorname{Im} \mu = 0$.
\end{theorem}

\begin{proof}
  \leanok
  Self-adjointness gives $\overline{\mu} = \mu$, and a complex number equal to
  its own conjugate has zero imaginary part. No finiteness is required.
\end{proof}

\begin{theorem}[Spectral theorem, matrix form]
  \label{thm:hermitian-spectral}
  \lean{QuantumBook.Foundations.hermitian_spectral_decomposition}
  \leanok
  \uses{thm:selfadjoint-real}
  A Hermitian matrix $A$ over $\mathbb{C}$ is diagonalized by a unitary change of
  basis: $A = U D U^{*}$, where $U$ is the unitary of an orthonormal eigenbasis
  and $D$ is the diagonal matrix of the (real) eigenvalues.
\end{theorem}

\begin{proof}
  \leanok
  \uses{thm:selfadjoint-real}
  Mathlib's \texttt{Matrix.IsHermitian.spectral\_theorem} phrases this with the
  star-algebra conjugation $\mathrm{conjStarAlgAut}$; unfolding it yields the
  explicit product $U D U^{*}$. The eigenvalues are real by
  Theorem~\ref{thm:selfadjoint-real}.
\end{proof}

\begin{theorem}[Spectral theorem, operator form]
  \label{thm:eigenvector-equation}
  \lean{QuantumBook.Foundations.selfAdjoint_eigenvector_equation}
  \leanok
  \uses{thm:hermitian-spectral}
  A self-adjoint operator $T$ on a finite-dimensional complex inner product space
  admits an orthonormal basis of eigenvectors $(e_i)$ with real eigenvalues
  $(\lambda_i)$, and acts on each by scaling: $T e_i = \lambda_i e_i$.
\end{theorem}

\begin{proof}
  \leanok
  \uses{thm:hermitian-spectral}
  Immediate from Mathlib's finite-dimensional diagonalization
  (\texttt{LinearMap.IsSymmetric.apply\_eigenvectorBasis}).
\end{proof}

\section{Frontier}

The finite-dimensional result above is the baby case of the general spectral
theorem. Its infinite-dimensional analogue, for unbounded self-adjoint operators,
is the first genuine frontier node of this project.

\begin{theorem}[Unbounded self-adjoint spectral theorem]
  \label{thm:unbounded-spectral}
  \notready
  \uses{thm:eigenvector-equation}
  An (unbounded) self-adjoint operator on a separable Hilbert space is unitarily
  equivalent to multiplication by a real measurable function on an
  $L^2$ space; equivalently, it admits a projection-valued spectral measure.
\end{theorem}

\begin{remark}
  This statement is \emph{not yet available in Mathlib} and carries no Lean
  declaration. It is tracked here as a frontier node (prompt.md invariant 2): the
  verified/frontier line runs exactly between
  Theorem~\ref{thm:eigenvector-equation} (proven) and
  Theorem~\ref{thm:unbounded-spectral} (frontier). Retiring it from the frontier,
  once Mathlib carries the unbounded spectral theorem, is a scheduled
  re-verification task (prompt.md \S11).
\end{remark}

\chapter{Part I -- Finite-dimensional quantum mechanics}

\section{The qubit and the Pauli operators}

The qubit's state space is $\mathbb{C}^2$; its traceless self-adjoint observables
are spanned by the Pauli operators $\sigma_x,\sigma_y,\sigma_z$. The three facts
below make them observables (Hermitian, hence real spectrum via
Theorem~\ref{thm:selfadjoint-real}), generators (traceless), and two-outcome
(involutions).

\begin{lemma}[Pauli operators are Hermitian]
  \label{lem:pauli-hermitian}
  \lean{QuantumBook.FiniteDim.sigmaZ_isHermitian}
  \leanok
  \uses{thm:selfadjoint-real}
  Each Pauli operator $\sigma_x,\sigma_y,\sigma_z$ is Hermitian; hence (by
  Theorem~\ref{thm:selfadjoint-real}) each has real eigenvalues and is a qubit
  observable. (Formalized for all three: \texttt{sigmaX/Y/Z\_isHermitian}.)
\end{lemma}
\begin{proof}\leanok Direct computation on the explicit matrices. \end{proof}

\begin{lemma}[Pauli operators are traceless]
  \label{lem:pauli-traceless}
  \lean{QuantumBook.FiniteDim.sigmaZ_trace}
  \leanok
  Each Pauli operator has trace $0$; the trace-zero self-adjoint $2\times2$
  matrices are the real span of the Paulis, i.e. $i\cdot\mathfrak{su}(2)$.
  (Formalized: \texttt{sigmaX/Y/Z\_trace}.)
\end{lemma}
\begin{proof}\leanok $\operatorname{tr} = A_{00}+A_{11}$; for $\sigma_z$ this is
  $1+(-1)=0$. \end{proof}

\begin{lemma}[Pauli operators are involutions]
  \label{lem:pauli-involution}
  \lean{QuantumBook.FiniteDim.sigmaZ_sq}
  \leanok
  \uses{thm:eigenvector-equation}
  Each Pauli operator squares to the identity, $\sigma^2 = I$. With the spectral
  theorem (Theorem~\ref{thm:eigenvector-equation}) this forces eigenvalues in
  $\{+1,-1\}$: a Pauli measurement has exactly two outcomes.
  (Formalized: \texttt{sigmaX/Y/Z\_sq}.)
\end{lemma}
\begin{proof}\leanok Direct $2\times2$ matrix multiplication. \end{proof}

\section{Observables and expectation values}

\begin{theorem}[Observables have real expectation values]
  \label{thm:expectation-real}
  \lean{QuantumBook.FiniteDim.expectation_isReal}
  \leanok
  \uses{thm:selfadjoint-real}
  If $A$ is a self-adjoint operator on a complex inner product space and $x$ is a
  state, then the expectation value $\langle A\rangle_x = \langle x, A x\rangle$ is
  real.
\end{theorem}
\begin{proof}\leanok
  Conjugating the inner product swaps its arguments ($\overline{\langle x, Ax\rangle}
  = \langle Ax, x\rangle$), and self-adjointness swaps them back
  ($\langle Ax, x\rangle = \langle x, Ax\rangle$); so $\langle A\rangle_x$ equals
  its own conjugate and is therefore real.
\end{proof}

\section{The Born rule}

\begin{theorem}[The Born rule is a probability distribution]
  \label{thm:born-rule}
  \lean{QuantumBook.FiniteDim.bornProb_sum_eq_one}
  \leanok
  \uses{thm:eigenvector-equation}
  Let $b$ be an orthonormal eigenbasis of an observable (Theorem
  \ref{thm:eigenvector-equation}) and $x$ a normalized state. The Born
  probabilities $p_i = |\langle b_i, x\rangle|^2$ are nonnegative and satisfy
  $\sum_i p_i = 1$; they form a genuine probability distribution over the
  outcomes.
\end{theorem}
\begin{proof}\leanok
  Nonnegativity is immediate. The sum equals $\|x\|^2$ by Parseval's identity for
  an orthonormal basis, which is $1$ on a normalized state.
\end{proof}

\section{Composite systems and entanglement}

\begin{theorem}[The Bell state is entangled]
  \label{thm:bell-entangled}
  \lean{QuantumBook.FiniteDim.bellState_isEntangled}
  \leanok
  The Bell state $b_0\otimes b_0 + b_1\otimes b_1$ on a two-level subsystem is not
  a product state: there are no $a, c$ with $b_0\otimes b_0 + b_1\otimes b_1 =
  a\otimes c$.
\end{theorem}
\begin{proof}\leanok
  In the tensor-product basis, a product $a\otimes c$ has coefficient
  $\alpha_i\gamma_j$ at $(i,j)$. Matching the Bell coefficients ($1$ on the
  diagonal, $0$ off) gives $\alpha_0\gamma_0 = 1$, $\alpha_0\gamma_1 = 0$,
  $\alpha_1\gamma_1 = 1$; the first forces $\alpha_0\ne 0$, so the second forces
  $\gamma_1 = 0$, contradicting the third.
\end{proof}

\section{Unitary evolution}

\begin{theorem}[Unitary evolution conserves total probability]
  \label{thm:probability-conservation}
  \lean{QuantumBook.FiniteDim.bornProb_sum_unitary}
  \leanok
  \uses{thm:born-rule}
  If $U$ is a unitary evolution (a linear isometry equivalence of the state space)
  then the total Born probability is invariant: $\sum_i p_i(Ux) = \sum_i p_i(x)$.
  A normalized state stays normalized under evolution.
\end{theorem}
\begin{proof}\leanok
  The total Born probability is the squared norm (Theorem \ref{thm:born-rule}), and
  a unitary preserves the norm.
\end{proof}

\begin{lemma}[Pauli algebra]
  \label{lem:pauli-algebra}
  \lean{QuantumBook.FiniteDim.sigmaX_mul_sigmaY}
  \leanok
  \uses{lem:pauli-involution}
  The Pauli operators satisfy $\sigma_x\sigma_y = i\sigma_z$ and its cyclic
  permutations. This is their multiplication table; it encodes the
  $\mathfrak{su}(2)$ Lie bracket and the Clifford relation generating qubit
  rotations. (Formalized: \texttt{sigmaX\_mul\_sigmaY},
  \texttt{sigmaY\_mul\_sigmaZ}, \texttt{sigmaZ\_mul\_sigmaX}.)
\end{lemma}
\begin{proof}\leanok Direct $2\times2$ matrix multiplication. \end{proof}
